\documentclass[11pt]{article}
\usepackage{amsmath, amssymb}
\usepackage{stmaryrd}          % for \llbracket and \rrbracket
\usepackage[english]{babel}
\usepackage{geometry}
\usepackage{tcolorbox}
\usepackage{booktabs}
\usepackage{array}
\geometry{margin=1in}

\title{PreExam Solution}
\author{}
\date{}

\begin{document}

\maketitle


\section{Ejercicio 1: Inferencia de Tipos}


Se analiza la función:

\begin{verbatim}
(define (count x l)
  (if (null? l)
      0
      (if (equal? (car l) x)
          (+ 1 (count x (cdr l)))
          (count x (cdr l)))))
\end{verbatim}

\subsection*{Paso 1: Identificación de variables}

\begin{itemize}
    \item $x$: tipo desconocido $\Rightarrow T$
    \item $l$: lista de elementos del mismo tipo que $x$
          $\Rightarrow T\ \mathbf{list}$
\end{itemize}

\subsection*{Paso 2: Análisis de operaciones}

\begin{itemize}
    \item \texttt{(null? l)} requiere $l : T\ \mathbf{list}$ y devuelve
          $\mathbf{bool}$
    \item \texttt{(car l)} devuelve el primer elemento $\Rightarrow T$
    \item \texttt{(cdr l)} devuelve la cola $\Rightarrow T\ \mathbf{list}$
    \item \texttt{equal?(car(l), x)} requiere ambos del mismo tipo
          $\Rightarrow$ ambos $T$, confirmando $x : T$
\end{itemize}

\subsection*{Paso 3: Tipo de retorno}

La función retorna:
\begin{itemize}
    \item $0$ en el caso base $\Rightarrow \mathbf{int}$
    \item $(+ \; 1 \; (\text{count} \ldots))$ en el caso recursivo
          $\Rightarrow \mathbf{int}$
\end{itemize}

Ambas ramas coinciden: el tipo de retorno es $\mathbf{int}$.

\subsection*{Resultado final}

\[
  \boxed{count : T \rightarrow T\ \mathbf{list} \rightarrow \mathbf{int}}
\]

\begin{tcolorbox}
  \textbf{Patrón clave:} funciones recursivas sobre listas suelen tener
  tipo $T \to T\ \mathbf{list} \to \alpha$.
  El tipo es \emph{polimórfico} en $T$: funciona para listas de enteros,
  strings, etc.
\end{tcolorbox}


\section{Ejercicio 2: Semántica Operacional}


Se desea definir la semántica de la instrucción:
\[
  \mathbf{repeat}\ E\ \mathbf{times}\ C
\]

\subsection*{Idea clave}

\begin{itemize}
    \item La expresión $E$ se evalúa \textbf{una sola vez} al inicio.
    \item El resultado $n$ es el contador; se decrementa en 1 después de
          cada ejecución de $C$.
    \item Introducimos una forma auxiliar
          $\mathbf{repeat\text{-}n}\ k\ \mathbf{times}\ C$
          donde $k \in \mathbb{N}$ ya es un valor entero (no una expresión).
\end{itemize}

% ── Big-step rules ────────────────────────────────────────────────────────────
\subsection*{Regla de entrada (evalúa $E$ una sola vez)}

\[
  \frac{
    \langle E,\, \sigma \rangle \Downarrow n
    \qquad
    \langle \mathbf{repeat\text{-}n}\ n\ \mathbf{times}\ C,\,
            \sigma \rangle \Downarrow \sigma'
  }{
    \langle \mathbf{repeat}\ E\ \mathbf{times}\ C,\,
            \sigma \rangle \Downarrow \sigma'
  }
  \quad \textsc{[Repeat-Eval]}
\]

\subsection*{Caso base ($k = 0$)}

\[
  \frac{\phantom{X}}{
    \langle \mathbf{repeat\text{-}n}\ 0\ \mathbf{times}\ C,\,
            \sigma \rangle \Downarrow \sigma
  }
  \quad \textsc{[Repeat-Zero]}
\]

\subsection*{Caso recursivo ($k > 0$)}


\[
  \frac{
    k > 0
    \qquad
    \langle C,\, \sigma \rangle \Downarrow \sigma'
    \qquad
    \langle \mathbf{repeat\text{-}n}\ (k{-}1)\ \mathbf{times}\ C,\,
            \sigma' \rangle \Downarrow \sigma''
  }{
    \langle \mathbf{repeat\text{-}n}\ k\ \mathbf{times}\ C,\,
            \sigma \rangle \Downarrow \sigma''
  }
  \quad \textsc{[Repeat-Step]}
\]

\begin{tcolorbox}
  \textbf{Corrección importante:} en la regla \textsc{Repeat-Step},
  \emph{no} se vuelve a evaluar $E$.
  Solo se trabaja con el contador entero $k$ y se resta $1$.
  Re-evaluar $E$ en cada iteración violaría la especificación del enunciado.\\[4pt]
  \textbf{Intuición:} equivalente a un \texttt{for}-loop donde el límite
  superior se fija antes de entrar al ciclo.
\end{tcolorbox}


\section{Ejercicio 3: Traducción Infix a Postfix}


\subsection*{Función semántica}

Definimos $\mathcal{P}\llbracket \cdot \rrbracket : \mathbf{Expr}
\to \mathbf{String}$ que traduce una expresión infija a postfija.

\subsection*{Reglas denotacionales}

\[
  \mathcal{P}\llbracket n \rrbracket = \texttt{"n"}
\]
\[
  \mathcal{P}\llbracket E_1 + E_2 \rrbracket
    = \mathcal{P}\llbracket E_1 \rrbracket
      \;\frown\; \mathcal{P}\llbracket E_2 \rrbracket
      \;\frown\; \texttt{"+"}
\]
\[
  \mathcal{P}\llbracket E_1 - E_2 \rrbracket
    = \mathcal{P}\llbracket E_1 \rrbracket
      \;\frown\; \mathcal{P}\llbracket E_2 \rrbracket
      \;\frown\; \texttt{"-"}
\]
\[
  \mathcal{P}\llbracket E_1 * E_2 \rrbracket
    = \mathcal{P}\llbracket E_1 \rrbracket
      \;\frown\; \mathcal{P}\llbracket E_2 \rrbracket
      \;\frown\; \texttt{"*"}
\]
\[
  \mathcal{P}\llbracket E_1 / E_2 \rrbracket
    = \mathcal{P}\llbracket E_1 \rrbracket
      \;\frown\; \mathcal{P}\llbracket E_2 \rrbracket
      \;\frown\; \texttt{"/"}
\]

donde $\frown$ denota concatenación de strings (con espacio separador).

\subsection*{Ejemplo: $3 + 4 * 5$}

Con precedencia estándar ($*$ liga más que $+$), el árbol es
$3 + (4 * 5)$:

\begin{align*}
  \mathcal{P}\llbracket 3 + (4 * 5) \rrbracket
    &= \mathcal{P}\llbracket 3 \rrbracket
       \;\frown\; \mathcal{P}\llbracket 4 * 5 \rrbracket
       \;\frown\; \texttt{"+"} \\
    &= \texttt{"3"}
       \;\frown\;
       \bigl(\mathcal{P}\llbracket 4 \rrbracket
             \;\frown\; \mathcal{P}\llbracket 5 \rrbracket
             \;\frown\; \texttt{"*"}\bigr)
       \;\frown\; \texttt{"+"} \\
    &= \texttt{"3"} \;\frown\; \texttt{"4"} \;\frown\;
       \texttt{"5"} \;\frown\; \texttt{"*"} \;\frown\; \texttt{"+"} \\
    &= \texttt{"3 4 5 * +"}
\end{align*}

\begin{tcolorbox}
  \textbf{Regla clave:} izquierdo $\frown$ derecho $\frown$ operador.
  La precedencia la maneja el \emph{parser} al construir el AST;
  la función $\mathcal{P}$ solo recorre el árbol resultante.
\end{tcolorbox}


\section{Ejercicio 4: Alcance Estático vs.\ Dinámico}


\subsection*{Los dos modelos}

\begin{center}
\begin{tabular}{lll}
  \toprule
  Modelo & Pregunta a hacer & Resultado \\
  \midrule
  \textbf{Estático (léxico)} &
    ¿Dónde está \emph{definido} \texttt{proc1} en el fuente? &
    $x = 2$ \\
  \textbf{Dinámico} &
    ¿Qué hay en la \emph{pila de llamadas} ahora mismo? &
    $x = 3$ \\
  \bottomrule
\end{tabular}
\end{center}

\subsection*{Alcance estático}

\texttt{proc1} está definido en el \textbf{nivel global} del programa.
El único $x$ visible en su contexto de definición es el $x$ global,
cuyo valor es $2$ (asignado por \texttt{main} antes de llamar a
\texttt{proc2}).

\[
  \Rightarrow \texttt{print} \text{ imprime } \textbf{x=2}
\]

\subsection*{Alcance dinámico}

En tiempo de ejecución, la pila al momento del \texttt{print} es:

\[
  \underbrace{\texttt{main}}_{\displaystyle x=2}
  \;\longrightarrow\;
  \underbrace{\texttt{proc2}}_{\displaystyle x=3 \text{ (local)}}
  \;\longrightarrow\;
  \texttt{proc1}
\]

El enlace más reciente de $x$ en la pila es el de \texttt{proc2},
con valor $3$.

\[
  \Rightarrow \texttt{print} \text{ imprime } \textbf{x=3}
\]

\begin{tcolorbox}
  \textbf{Regla mental:}\\
  Estático $=$ dónde se \emph{define} (contexto léxico)\\
  Dinámico $=$ desde dónde se \emph{llama} (pila en tiempo de ejecución)
\end{tcolorbox}


\section{Ejercicio 5: Paso de Parámetros}


\subsection*{Tipos de parámetros}

\begin{itemize}
    \item \textbf{ref} (por referencia): alias directo al argumento.
          Leer/escribir el parámetro lee/escribe la variable original.
    \item \textbf{name} (por nombre / Jensen's device):
          la expresión argumento se \emph{re-evalúa} cada vez que se usa
          el parámetro.  Si las variables internas cambian, cada uso puede
          acceder a una celda diferente.
    \item \textbf{VR} (value-result): copia \emph{in} al entrar
          ($z_{\text{local}} \leftarrow $ valor actual del argumento) y
          copia \emph{out} al salir ($i \leftarrow z_{\text{local}}$).

\end{itemize}

\subsection*{Estado inicial}

\[
  a = \{0,\,10,\,20,\,40\}, \qquad i = 0
\]

\subsection*{Llamada: \texttt{test(i, a[i], i)}}

\begin{itemize}
  \item $x \xleftarrow{\text{ref}} i$: alias de $i$
  \item $y \xleftarrow{\text{name}} \mathtt{a[i]}$: expresión diferida
  \item $z_{\text{local}} \xleftarrow{\text{VR}} 0$ (copia-in de $i=0$)
\end{itemize}

\subsection*{Ejecución línea a línea}

\begin{center}
\renewcommand{\arraystretch}{1.4}
\begin{tabular}{clll}
  \toprule
  Línea & Sentencia & Efecto & Estado tras ejecutar \\
  \midrule
  8 & \texttt{x := x+1} &
      $x$ es alias de $i$; $i \leftarrow 0{+}1 = 1$ &
      $i=1,\; z_{\text{local}}=0$ \\
  9 & \texttt{y := y+1} &
      $y = \mathtt{a[i]}$; $i=1$ ahora $\Rightarrow y = a[1] = 10$ &
      $a[1] \leftarrow 11$ \\
  10 & \texttt{z := z+2} &
      copia local $z_{\text{local}} \leftarrow 0{+}2 = 2$ &
      $z_{\text{local}}=2$, $i=1$ \\
  11 & \texttt{y := y+1} &
      $y = \mathtt{a[i]}$; $i=1$ todavía $\Rightarrow y = a[1] = 11$ &
      $a[1] \leftarrow 12$ \\
  \midrule
  \multicolumn{2}{l}{\textbf{Retorno (VR copia-out)}} &
      $z_{\text{local}} = 2 \;\Rightarrow\; i \leftarrow 2$ &
      $i = 2$ \\
  \bottomrule
\end{tabular}
\end{center}

\subsection*{Resultado final}


\[
  \boxed{i = 2, \qquad a = \{0,\; 12,\; 20,\; 40\}}
\]

\begin{tcolorbox}
  \textbf{Tres interacciones críticas:}
  \begin{enumerate}
    \item \textbf{ref cambia el índice de name.}
          La línea 8 incrementa $i$ de 0 a 1 vía \texttt{ref}.
          Las líneas 9 y 11 usan $y = a[i]$ con $i=1$, accediendo a
          $a[1]$ (no $a[0]$).
    \item \textbf{name modifica la misma celda dos veces.}
          Tanto la línea 9 como la 11 operan sobre $a[1]$,
          incrementándolo de 10 a 11 y luego a 12.
    \item \textbf{VR sobreescribe el cambio de ref al retornar.}
          \texttt{ref} dejó $i=1$, pero la copia-out de VR escribe
          $z_{\text{local}}=2$ sobre $i$, resultando en $i=2$.
  \end{enumerate}
\end{tcolorbox}

\subsection{Worked Examples}

The following examples are designed to address two common sources of
confusion: the difference between big-step and small-step style, and
how to extend the semantics with new constructs. Each example is worked
out in full.

% ---------------------------------------------------------------
\subsubsection{Example 1: Big-step vs.\ Small-step on the Same Program}
% ---------------------------------------------------------------

Consider the following SIMP program with initial state
$\sigma = \{x \mapsto 3,\; y \mapsto 0\}$:

\begin{verbatim}
y := !x + 1 ;
x := !y * 2
\end{verbatim}

\paragraph{Big-step derivation.}
Big-step semantics derives the \emph{final} state in one tree. We write
$\langle C, \sigma \rangle \Downarrow \sigma'$ to mean ``command $C$
in state $\sigma$ terminates in state $\sigma'$''.

The key insight: in big-step style you build a \textbf{derivation tree} the root is the whole program, the leaves are axioms. Work
\emph{top-down} by applying the sequence rule, then fill in each
sub-derivation.

\[
\dfrac{
  \dfrac{
    \langle !x + 1,\;\sigma \rangle \Downarrow 4
  }{
    \langle y := !x + 1,\;\sigma \rangle \Downarrow \sigma_1
  }
  \qquad
  \dfrac{
    \langle !y * 2,\;\sigma_1 \rangle \Downarrow 8
  }{
    \langle x := !y * 2,\;\sigma_1 \rangle \Downarrow \sigma_2
  }
}{
  \langle y := !x + 1\;;\; x := !y * 2,\;\sigma \rangle
  \Downarrow \sigma_2
}
\]

where $\sigma_1 = \sigma[y \mapsto 4]$ and
$\sigma_2 = \sigma_1[x \mapsto 8]$.

Final state: $\sigma_2 = \{x \mapsto 8,\; y \mapsto 4\}$.

\paragraph{Small-step derivation.}
Small-step semantics reduces the program \emph{one step at a time}.
We write $\langle C, \sigma \rangle \to \langle C', \sigma' \rangle$
(or $\langle C, \sigma \rangle \to \langle \mathbf{skip}, \sigma'
\rangle$ at the last step).

The key insight: in small-step style you write a \textbf{linear
sequence} of configurations. At each step, apply exactly one rule to
the leftmost reducible sub-expression.

\begin{align*}
&\langle\; y := !x + 1\;;\; x := !y * 2,\;\sigma \rangle \\
\to\;&\langle\; y := 3 + 1\;;\; x := !y * 2,\;\sigma \rangle
  &&\text{(dereference $!x$)} \\
\to\;&\langle\; y := 4\;;\; x := !y * 2,\;\sigma \rangle
  &&\text{(arithmetic)} \\
\to\;&\langle\; \mathbf{skip}\;;\; x := !y * 2,\;\sigma_1 \rangle
  &&\text{(assign, $\sigma_1 = \sigma[y\mapsto 4]$)} \\
\to\;&\langle\; x := !y * 2,\;\sigma_1 \rangle
  &&\text{(skip-seq)} \\
\to\;&\langle\; x := 4 * 2,\;\sigma_1 \rangle
  &&\text{(dereference $!y$)} \\
\to\;&\langle\; x := 8,\;\sigma_1 \rangle
  &&\text{(arithmetic)} \\
\to\;&\langle\; \mathbf{skip},\;\sigma_2 \rangle
  &&\text{(assign, $\sigma_2 = \sigma_1[x\mapsto 8]$)}
\end{align*}

Both styles reach the same final state
$\{x \mapsto 8,\; y \mapsto 4\}$.

\paragraph{When to use each.}
\begin{itemize}
  \item Use \textbf{big-step} when you only care about the final result
  and want a compact derivation.
  \item Use \textbf{small-step} when you need to reason about
  intermediate states, non-termination, or concurrency.
  \item Never mix styles in the same derivation.
\end{itemize}

\subsubsection{Example 2: Extending SIMP with \texttt{for}}

We extend SIMP with the construct \texttt{for} $E_1$ \texttt{to} $E_2$
\texttt{do} $C$, which evaluates $E_1$ and $E_2$ \emph{once} at the
start and executes $C$ for each integer $i$ with $n_1 \leq i \leq n_2$
(does nothing if $n_1 > n_2$).

\paragraph{Step 1 Define the rules}

The standard strategy is to \emph{desugar} the new construct into
existing ones once the bounds are known. We need three rules.

\medskip
\noindent\textbf{Evaluate the bounds} (two expression-reduction steps,
small-step):
\[
  \frac{
    \langle E_1, \sigma \rangle \to \langle E_1', \sigma \rangle
  }{
    \langle \texttt{for}\; E_1\;\texttt{to}\; E_2\;\texttt{do}\; C,\;
    \sigma \rangle \;\to\;
    \langle \texttt{for}\; E_1'\;\texttt{to}\; E_2\;\texttt{do}\; C,\;
    \sigma \rangle
  }
\]
\[
  \frac{
    \langle E_2, \sigma \rangle \to \langle E_2', \sigma \rangle
  }{
    \langle \texttt{for}\; n_1\;\texttt{to}\; E_2\;\texttt{do}\; C,\;
    \sigma \rangle \;\to\;
    \langle \texttt{for}\; n_1\;\texttt{to}\; E_2'\;\texttt{do}\; C,\;
    \sigma \rangle
  }
\]

\noindent\textbf{Base case} bounds are values, loop is empty
($n_1 > n_2$):
\[
  \frac{n_1 > n_2}{
    \langle \texttt{for}\; n_1\;\texttt{to}\; n_2\;\texttt{do}\; C,\;
    \sigma \rangle \;\to\;
    \langle \mathbf{skip},\; \sigma \rangle
  }
\]

\noindent\textbf{Inductive case} unfold one iteration ($n_1 \leq
n_2$):
\[
  \frac{n_1 \leq n_2}{
    \langle \texttt{for}\; n_1\;\texttt{to}\; n_2\;\texttt{do}\; C,\;
    \sigma \rangle \;\to\;
    \langle C\;;\;\texttt{for}\; (n_1{+}1)\;\texttt{to}\; n_2\;
    \texttt{do}\; C,\; \sigma \rangle
  }
\]

\paragraph{Step 2 Trace through an example.}

Use state $\sigma = \{x \mapsto 0,\; n_1 \mapsto 1,\; n_2 \mapsto 3\}$
and program \texttt{for} $!n_1$ \texttt{to} $!n_2$ \texttt{do}
$x := !x + 1$.

\begin{align*}
&\langle \texttt{for}\;!n_1\;\texttt{to}\;!n_2\;\texttt{do}\;
  x{:=}!x{+}1,\;\sigma \rangle \\
\to\;&\langle \texttt{for}\;1\;\texttt{to}\;!n_2\;\texttt{do}\;
  x{:=}!x{+}1,\;\sigma \rangle
  &&\text{(deref $!n_1$)}\\
\to\;&\langle \texttt{for}\;1\;\texttt{to}\;3\;\texttt{do}\;
  x{:=}!x{+}1,\;\sigma \rangle
  &&\text{(deref $!n_2$)}\\
\to\;&\langle x{:=}!x{+}1\;;\;\texttt{for}\;2\;\texttt{to}\;3\;
  \texttt{do}\;x{:=}!x{+}1,\;\sigma \rangle
  &&\text{(unfold, $1 \leq 3$)}\\
\to\;&\langle \mathbf{skip}\;;\;\texttt{for}\;2\;\texttt{to}\;3\;
  \texttt{do}\;x{:=}!x{+}1,\;\sigma_1 \rangle
  &&\text{(assign, $\sigma_1 = \sigma[x\mapsto 1]$)}\\
\to\;&\langle \texttt{for}\;2\;\texttt{to}\;3\;\texttt{do}\;
  x{:=}!x{+}1,\;\sigma_1 \rangle
  &&\text{(skip-seq)}\\
  &\quad\vdots \\
\to\;&\langle \mathbf{skip},\;\sigma_3 \rangle
  &&\text{(base case, $4 > 3$, $\sigma_3 = \sigma[x \mapsto 3]$)}
\end{align*}

\paragraph{Common mistakes to avoid.}
\begin{itemize}
  \item \textbf{Evaluating the bounds repeatedly.} The bounds $n_1$ and
  $n_2$ must be fully reduced to values \emph{before} unfolding begins.
  If you unfold first, you may re-evaluate a side-effectful expression
  on every iteration.
  \item \textbf{Forgetting the base case.} Without the $n_1 > n_2$
  rule, the derivation never terminates for empty ranges.
  \item \textbf{Mixing big-step and small-step notation.} In big-step,
  the \texttt{for} rules look different, each premise is a full
  $\Downarrow$ judgment, not a single $\to$ step.
\end{itemize}

\end{document}